%Notes from Cliff Toubes, transcribed with ChatGPT
%\documentclass[12pt,letterpaper]{article}
%\usepackage{amsmath, amssymb, amsthm}


%%%
%% Required packages:
%%
\usepackage{etex}
\usepackage{amsmath}
\usepackage{amssymb}
\usepackage{amstext}
\usepackage{amsthm}
\usepackage{fancyhdr,fancybox}
\usepackage{mathrsfs} % need for \mathscr command
\usepackage{multicol}
\usepackage[normalem]{ulem}
%\usepackage{pictex,pictexplus}
% \usepackage{pictexwd}
\usepackage{rotating}
\usepackage{url}
\usepackage{xfrac}
\usepackage{booktabs}
\usepackage{color,soul}
\usepackage{comment}

\usepackage{hyperref}
\hypersetup{
    colorlinks=true,     % false: boxed links; true: colored links
    linkcolor=blue,      % color of internal links
    citecolor=blue,      % color of links to bibliography
    filecolor=blue,      % color of file links
    urlcolor=blue        % color of external links
}


\usepackage{tikz}
\usetikzlibrary{backgrounds,calc}

%%
%% Common formatting commands
%%
\setlength{\parindent}{0in}
\setlength{\textwidth}{7in}
\setlength{\evensidemargin}{-0.25in}
\setlength{\oddsidemargin}{-0.25in}
\setlength{\parskip}{.5\baselineskip}
\setlength{\topmargin}{-0.5in}
\setlength{\textheight}{9in}

%               Problem and Part
\newcounter{problemnumber}
\newcounter{partnumber}
\newcounter{subpartnumber}
\newcommand{\Problem}{%
  \refstepcounter{problemnumber}%
  \setcounter{partnumber}{0}%
  \setcounter{subpartnumber}{0}%
  \item[\makebox{\hfil\textbf{\theproblemnumber.}\hfil}]%
  }
\newcommand{\InProblem}[2][1.6]{%
  \refstepcounter{problemnumber}%
  \setcounter{partnumber}{0}%
  \setcounter{subpartnumber}{0}%
  \textbf{\theproblemnumber.}\ \ \parbox[t]{#1in}{#2}%
}
\newcommand{\InSmallProblem}[1]{\InProblem[1.05]{#1}}
\newcommand{\NoProblem}[1]{%
  \mbox{\hfil\textbf{#1}\hfil}%
  }
\newcommand{\UnProblem}{%
  \refstepcounter{problemnumber}%
  \setcounter{partnumber}{0}%
  \setcounter{subpartnumber}{0}%
  \mbox{\hfil\textbf{\theproblemnumber.}\hfil}%
  }
\newcommand{\InFlexProblem}[1]{%
  \refstepcounter{problemnumber}%
  \setcounter{partnumber}{0}%
  \setcounter{subpartnumber}{0}%
  \textbf{\theproblemnumber.}\ \ #1%
}
%%  Part:
\newcommand{\Part}{%
  \renewcommand{\thepartnumber}{(\alph{partnumber})}%
  \refstepcounter{partnumber}
  \setcounter{subpartnumber}{0}%
  \item[(\alph{partnumber})]%
}
\newcommand{\InPart}[2][1.75]{%
  \renewcommand{\thepartnumber}{(\alph{partnumber})}%
  \refstepcounter{partnumber}%
  \setcounter{subpartnumber}{0}%
  (\alph{partnumber})\ \ \parbox[t]{#1in}{#2}%
}
\newcommand{\UnPart}{%
  \refstepcounter{partnumber}%
  \renewcommand{\thepartnumber}{(\alph{partnumber})}%
  \setcounter{subpartnumber}{0}%
  (\alph{partnumber})%
}
\newcommand{\InLargePart}[1]{\InPart[3.05]{#1}}
\newcommand{\InFullPart}[1]{\InPart[6.2]{#1}}
\newcommand{\InSmallPart}[1]{\InPart[1.05]{#1}}
%% SubPart:
\newcommand{\SubPart}{%
  \refstepcounter{subpartnumber}%
  \item[(\roman{subpartnumber})]%
  }
\newcommand{\InSubPart}[2][2.25]{%
  \stepcounter{subpartnumber}%
  (\roman{subpartnumber})\ \ \parbox[t]{#1in}{#2}%
}
\newcommand{\InSmallSubPart}[1]{\InSubPart[1.5]{#1}}
\newcommand{\InTinySubPart}[1]{%
  \stepcounter{subpartnumber}%
  (\roman{subpartnumber})\ \ \makebox{#1}%
}
\newcommand{\InMySubPart}[2][2.25]{%
  \stepcounter{subpartnumber}%
  \parbox[t]{#1in}{\makebox[0.3in][l]{(\roman{subpartnumber})}\ \ {#2}}%
}
\newcommand{\TrueFalse}{\textbf{T}\ \ \textbf{F}\ \ }
\newcommand{\TFPart}[1]{% 
  \stepcounter{partnumber}%
  \setcounter{subpartnumber}{0}%
  \item[(\alph{partnumber})] \TrueFalse\parbox[t]{5.5in}{#1}}

\newcommand{\Points}[1]{(#1 points) \ }
\newcommand{\Point}[1]{(#1 point) \ }


%% For Answers:
\newcounter{answernumber}
\newcommand{\Answer}{%
  \refstepcounter{answernumber}%
  \setcounter{partnumber}{0}%
  \setcounter{subpartnumber}{0}%
  \item[\makebox{\hfil\textbf{\theanswernumber.}\hfil}]%
  }


% Setting up the headers for the first page
%\chead{} % will default to what is typed in the file.
\fancyhead[L]{\textsc{Math 22a}}
\fancyhead[R]{\textsc{Fall 2024}}
\fancyfoot[R]{
	\vspace*{\fill}\footnotesize\textcopyright{} \emph{Harvard Math 22a}	
}
\renewcommand{\headrulewidth}{0pt}

%\fancyhead[C]{}
%	\chead{}	
%	\lhead{}
%	\rhead{}
%	\cfoot{}


% Putting copyright on the footer of each page
%\fancypagestyle{secondpage}{
%	\fancyhead[C]{TESTing again} % change this to be blank
%		%\chead{}	
%	\fancyhead[L]{bears} % change this to be blank
%	\fancyhead[R]{dogs} % change this to be blank
%	\fancyfoot[R]{ %keeps the footer the same
%		\vspace*{\fill}\footnotesize\textcopyright{} \emph{Harvard Math 22a}	
%	}
%}
%\renewcommand{\headrulewidth}{0pt}
%\pagestyle{fancy}
\pagestyle{empty}


%\lhead{}
%\rhead{}
%\rfoot{\vspace*{\fill}\footnotesize\textcopyright{} \emph{Harvard Math 22a}}
%\renewcommand{\headrulewidth}{0pt}
%\pagestyle{fancy}




%%
%% Common shortcut commands:
%%
\newcommand{\ds}{\displaystyle}
\newcommand{\sgn}{\operatorname{sgn}}
\renewcommand{\Pr}{\operatorname{P}}
\newcommand{\CPr}[3][{}]{\Pr_{\text{#1}}\left(#2\,|\,#3\right)}

\newcommand{\C}{\mathbb{C}}
\newcommand{\R}{\mathbb{R}}
\newcommand{\Z}{\mathbb{Z}}
\newcommand{\N}{\mathbb{N}}
\newcommand{\Q}{\mathbb{Q}}
\newcommand{\D}{\mathbb{D}}

\newcommand{\rref}{\operatorname{rref}}
\newcommand{\im}{\operatorname{im}}
\newcommand{\rank}{\operatorname{rank}}
\newcommand{\diag}{\operatorname{diag}}
\newcommand{\Span}{\operatorname{span}}
%\renewcommand{\vec}[1]{\mathbf{#1}}
\newcommand{\charpoly}{\mathscr{P}}
\newcommand{\nicefrac}[2]{\sfrac{#1}{#2}} % using xfrac package
\newcommand{\topology}{\mathcal{T}}
\newcommand{\Inn}{\operatorname{Inn}}

\newcommand{\blankbox}[2]{\fbox{\rule{#1}{0in}\rule{0in}{#2}}}

\newenvironment{myindentpar}[1]%
 {\begin{list}{}%
         {\setlength{\leftmargin}{#1}
          \setlength{\rightmargin}{\leftmargin}}%
         \item[]%
 }
 {\end{list}}
\newtheorem{theorem}{Theorem}
\newtheorem*{NStheorem}{Theorem}
\newtheorem{cor}[theorem]{Corollary}
\newtheorem*{claim}{Claim}
\newtheorem{defn}{Definition}

\newenvironment{amatrix}[1]{%
  \left(\begin{array}{@{}*{#1}{c}|c@{}}
}{%
  \end{array}\right)
}

%--------grstep
% For denoting a Gauss' reduction step.
% Use as: \grstep{\rho_1+\rho_3} or \grstep[2\rho_5 \\ 3\rho_6]{\rho_1+\rho_3}
\newcommand{\grstep}[2][\relax]{%
   \ensuremath{\mathrel{
       {\mathop{\longrightarrow}\limits^{#2\mathstrut}_{
                                     \begin{subarray}{l} #1 \end{subarray}}}}}}
\newcommand{\swap}{\leftrightarrow}



\section{{Singular Value Decomposition, $\vec\S$7.5}}

\chead{\Large\textbf{Singular Value Decomposition, $\vec\S$7.5}}

%\begin{document}
\thispagestyle{fancy}

\textbf{Review from previously:}
    \begin{enumerate}
    \setlength{\itemsep}{3em}
        \item \textbf{Definition:} A quadratic form on $\mathbb{R}^n$ is a function $Q$ whose value at $x$ can be computed using the rule $Q(x) = x^T A x = x \cdot A x$ where $A$ is a symmetric matrix.
        \item Minkowski 'inner product'. Write $\mathbb{R}^{n+1}$ as $(t, y)$ with $y \in \mathbb{R}^n$.
        \begin{enumerate}
            \item $Q((t, y)) = \|y\|^2 - t^2$.
            \item The theory of special relativity is all about this Minkowski inner product.
        \end{enumerate}
        \item Change of variable
        \begin{enumerate}
            \item If $x = Py$ where $P$ is invertible, then $x^T A x = y^T P^T A P y$.
            \item Thus, in the new coordinates, $Q$ is represented not by $A$, but by $P^T A P$.
        \end{enumerate}
        \item Principle axis theorem: There is a change of variable matrix $P$ that is orthonormal such that $Q(y) = y^T D y$ with $D$ being diagonal (the eigenvalues of $A$ are its diagonal entries).
        \item Constrained optimization: Look at $Q(x)$ subject to the constraint that $\|x\| = 1$.
        \begin{enumerate}
            \item If $A$ is the symmetric matrix such that $Q(x) = x^T A x$, then the maximum of $Q$ on the constraint surface is the largest eigenvalue of $A$; and likewise the minimum is the smallest eigenvalue of $A$.
            \item Minimize $Q$ on the $\|x\| = 1$ sphere subject to the constraint that $x$ be orthogonal to all eigenvectors previously found.
        \end{enumerate}
       \end{enumerate}
\newpage

\begin{enumerate}
    \setlength{\itemsep}{5em}

    \item We looked at symmetric matrices: Let $A$ be a symmetric $n \times n$ matrix.
    \begin{enumerate}    \setlength{\itemsep}{3em}

        \item The matrix $A$ is orthogonally diagonalizable: $A = P D P^{-1}$ with $P^{-1} = P^T$.
        \item The columns of $P$ are an orthonormal basis $\{v_1, \dots, v_n\}$ of eigenvectors for $A$.
        \item The eigenvalues of $A$ are real and are the diagonal entries of $D$.
    \end{enumerate}

    \item What can be said about an $m \times n$ matrix $A$ when $m \neq n$?
    \begin{enumerate}    \setlength{\itemsep}{3em}

        \item $A^T A$ is a symmetric, $n \times n$ matrix. (We used this matrix $A^T A$ when doing least squares and orthogonal projection to a subspace in $\mathbb{R}^n$.)
        \item Thus, it is orthogonally diagonalizable: $A^T A = P D P^{-1}$ with $P^{-1} = P^T$.
        \item The eigenvalues are non-negative because $v \cdot A^T A v = (A v) \cdot (A v) = \|A v\|^2$.
        \item The square roots of the eigenvalues $\sigma_1 = \sqrt{\lambda_1}, \dots, \sigma_n = \sqrt{\lambda_n}$ are called the singular values.
        \item $\text{Null}(A)$ has dimension at least $n - m$, so at most $m$ of the $\lambda_k$'s are not zero.
    \end{enumerate}

\newpage

    \item We can likewise look at $A A^T$ which is an $m \times m$ symmetric matrix.
    \begin{enumerate}    \setlength{\itemsep}{5em}

        \item It can be written as $Q B Q^{-1}$ where $Q$ is an $m \times m$ orthogonal matrix and $B$ is diagonal with entries being the eigenvalues of $A A^T$.
    \end{enumerate}

    \item \textbf{Theorem:} The non-zero eigenvalues of $A A^T$ are the same as those of $A^T A$.
    \begin{enumerate}    \setlength{\itemsep}{5em}

        \item Indeed, if $A^T A v = \lambda v$, then $A (A^T A v) = \lambda A v$ and $A (A^T A v)$ is $(A A^T) A v$.
        \item Also, if $A A^T u = \lambda' u$, then $A^T (A A^T u) = \lambda' A^T u$.
        \item So if $v_k$ is an eigenvector of $A^T A$ with non-zero eigenvalue $\lambda_k$ and length 1, then $u_k = \lambda_k^{-1/2} A v_k$ is an eigenvector of $A A^T$ with eigenvalue $\lambda_k$ and length 1.
        \item Moreover $u_k \cdot u_j = \lambda_k^{-1/2} \lambda_j^{-1/2} (A v_k) \cdot (A v_j) = \lambda_k^{-1/2} \lambda_j^{1/2} v_k \cdot A^T A v_j = \lambda_k^{-1/2} \lambda_j^{1/2} v_k \cdot v_j$.
        \item Thus, the $\{ u_k = \sigma_k^{-1} A v_k : \lambda_k \neq 0\}$ are an orthonormal basis for the eigenspaces of $A A^T$ with non-zero eigenvalue.
    \end{enumerate}

\newpage     \setlength{\itemsep}{7em}


    \item Thus, $A A^T$ and $A^T A$ have the same set of non-zero eigenvalues (which are non-negative). List these as $\lambda_1, \dots, \lambda_r$ for $r \leq \min(m, n)$.
    \item We can then write $A = \sigma_1 u_1 v_1^T + \cdots + \sigma_r u_r v_r^T$ where $\{v_i\}_{i=1,\dots,r}$ are an orthonormal basis for the eigenspaces of $A$ with non-zero eigenvalue, and $u_k = \sigma_k^{-1} A v_k$.
    \begin{enumerate}    \setlength{\itemsep}{5em}

        \item Indeed, if $x$ is any vector in $\mathbb{R}^n$, I can use the basis $\{v_1, \dots, v_r, v_{r+1}, \dots, v_n\}$ to write $x$ as $x = c_1 v_1 + \cdots + c_n v_n$.
        \item Then $A x = \sigma_1 c_1 u_1 + \cdots + \sigma_r c_r u_r$ which gives the expansion of $A x$ in terms of the $u$-basis. (Note that if $A A^T u = 0$, then $A^T u = 0$ so $A x$ has no $u_k$'s in its expansion from the 0 eigenvalues of $A A^T$.) Likewise, there are no $c_k$ for those $v_k$ with $A v_k$ being 0.
        \item In this regard, $u_k \cdot A x = (A^T u_k) \cdot x = \sigma_k v_k \cdot x = \sigma_k c_k$ as it should be.
    \end{enumerate}

    \item The depiction of $A$ as $\sigma_1 u_1 v_1^T + \cdots + \sigma_r u_r v_r^T$ is called the singular value decomposition of $A$.

\end{enumerate}
\newpage

\chead{\Large\textbf{Arc of Math 22A}}
%\thispagestyle{fancy}

 %\textbf{\Large Arc of Math 22A:}
 
 \begin{enumerate}
        \item We started with linear equations, wrote them in terms of matrices ($A x = b$) and then introduced the row-reduction algorithm to solve them or find them inconsistent by obtaining $\text{rref}(A)$.
        \item We then started playing with the notion of a matrix:
        \begin{enumerate}
            \item Multiplying them against vectors, multiplying them against other matrices.
            \item The notion of an invertible matrix and its inverse.
            \item The set of invertible matrices form a group denoted by $Gl(n; \mathbb{R})$.
            \item Introduced elementary matrices to formalize the row reduction operation.
        \end{enumerate}
        \item

 Then, we discovered the determinant of a matrix which gives us a way to determine when a matrix is invertible.
        \begin{enumerate}
            \item Lots of neat properties of determinants ($\det(AB) = \det(A) \det(B)$).
            \item An inductive formula for the determinant using expansion across any row or column.
        \end{enumerate}
        \item Introduced the notion of a vector space:
        \begin{enumerate}
            \item Subspaces in $\mathbb{R}^n$.
            \item Null space and column space of a matrix.
            \item The notion of a basis for a vector space giving coordinates for the vector space, which allows us to write linear transformations as matrix multiplication.
            \item Change of basis formulas.
            \item Notion of dimension of a subspace and rank of a matrix (dimension of its column space).
        \end{enumerate}
        \item The notion of an eigenvector and eigenvalue for a square matrix:
        \begin{enumerate}
            \item Characteristic equation for potential eigenvalues.
            \item Diagonalizable matrices (not all are diagonalizable over $\mathbb{R}$).
            \item Complex matrices and complex eigenvalues (most matrices are diagonalizable if complex eigenvalues are allowed).
            \item Using diagonalization to solve discrete dynamical systems and linear differential equations.
        \end{enumerate}
        \item Introduced dot products (inner products) and studied the many things it brings with it:
        \begin{enumerate}
            \item Notion of length, distance, orthogonality.
            \item Orthogonal basis, orthonormal basis, orthogonal projections.
            \item Gram-Schmidt algorithm.
            \item Least squares and finding the closest vector in a subspace.
        \end{enumerate}
        \item Continuing with dot products, we studied symmetric matrices:
        \begin{enumerate}
            \item These are orthogonally diagonalizable (and vice versa).
            \item There is an orthonormal basis consisting of eigenvectors. Eigenvalues are real.
            \item Spectral theorem.
            \item Singular value decomposition for $m \times n$ matrices even when $m \neq n$.
        \end{enumerate}
    \end{enumerate}


     \textbf{And} we learned how to write proofs and interpret rigorous mathematics.
    
     \textbf{And} we learned something about science and data.

%TODO
% Is there anything to do in the way of solutions here? Or is it just an outline of notes?  Fine if that is what you end up with, but put a blank page.  Or maybe just ask the students to answer the question "Why?" after each declarative statement.


%\end{document}
